% =====================================================================
% Chapter 2 revisions for the final pipeline (EQUITAS-RCMF on UBFC-Phys)
% Replace the corresponding passages in chapters/chapter_2.tex.
% Only the "Relevance to this thesis" paragraphs (and one 2.3 paragraph)
% change; the summaries and limitations of each paper stay as they are.
% =====================================================================


% ---------------------------------------------------------------------
% 1. Schmidt et al. (2018), WESAD -- replace the "Relevance" paragraph
% ---------------------------------------------------------------------
\textit{Relevance to this thesis.} Although the final pipeline uses UBFC-Phys rather than WESAD, this benchmark informs it in three ways. First, both datasets induce stress with the Trier Social Stress Test, so the stress label in this thesis is the same kind of controlled laboratory proxy, not a direct measure of threat-related states. Second, the benchmark found wrist-worn signals less informative than chest-worn ones (88\% versus 93.12\% for stress versus non-stress), and the physiological input in this thesis comes from a wrist-worn device, the weaker of the two sensor locations studied. Third, the benchmark's leave-one-subject-out evaluation shows how results on a small cohort should be assessed: with eight participants, performance on unseen individuals can only be estimated if each participant's data are kept entirely in either training or testing.


% ---------------------------------------------------------------------
% 2. Costantini et al. (2023) -- replace the "Relevance" paragraph
% ---------------------------------------------------------------------
% TODO(authors): Confirm in Sabour et al. (2023) that UBFC-Phys records BVP and EDA
% with the Empatica E4, and check which BVP/EDA features the model uses (e.g. mean IBI,
% SDNN, RMSSD, EDA statistics) so the second-to-last sentence matches the code.
\textit{Relevance to this thesis.} Costantini et al.\ validated the Empatica E4, the wrist device with which UBFC-Phys records the BVP and EDA signals used in this thesis, so their findings bear directly on the quality of the model's physiological input. Without processing, the device captured only about half of the inter-beat intervals; wrist EDA showed no agreement with the reference in any condition; and among HRV features, the mean inter-beat interval remained reliable while RMSSD lost its agreement under movement. These limits matter more than usual for this model. Because the fused model shows modality collapse, relying on physiology more than on vision, the reliability of the physiological signals directly bounds its performance. Features derived from BVP should therefore favour the more robust measures, such as the mean inter-beat interval, and conclusions drawn from EDA-derived features should be treated with caution. The findings also support the sensor-noise robustness evaluation proposed in Section~1.6.


% ---------------------------------------------------------------------
% 3. Toneva et al. (2019) -- replace the "Relevance" paragraph
% ---------------------------------------------------------------------
\textit{Relevance to this thesis.} Toneva et al.\ provide a per-example measure of learning difficulty that can be computed during training. Applied to this thesis, forgetting counts could show whether scar-positive windows are learned less stably than scar-negative ones, complementing the post-compression analysis suggested by Hooker et al.~\cite{hooker2019}: one identifies the examples a model finds hard to learn, the other those it loses when compressed. With only eight participants, each contributing many correlated windows, forgetting statistics aggregated by participant would also show whether particular individuals' data are consistently difficult. Their finding that mislabelled examples are the most forgotten is relevant to how stress is labelled here. Every window from the speech and arithmetic tasks is labelled as stress and every rest window as non-stress, although individual physiological responses vary within a task; windows whose physiology does not reflect their task label behave, from the model's perspective, like noisy labels.


% ---------------------------------------------------------------------
% 4. Vapnik and Vashist (2009) -- replace the "Relevance" paragraph
% ---------------------------------------------------------------------
\textit{Relevance to this thesis.} LUPI formalizes the setting of this thesis: the scar segmentation mask is available during training but not at deployment~\cite{vapnik2009, vapnik2015}. In EQUITAS-RCMF, the mask is used during training to compute the focus signal that measures how strongly visual features concentrate on the scar, while at inference the focus is estimated from the confounder subspace and the mask pathway is absent from the compiled model. The two works use training-only information for different purposes. LUPI uses it to estimate the difficulty of each training example and so improve accuracy, whereas this thesis uses it to route scar information away from the causal representation and to suppress visual evidence that concentrates on the scar. Vapnik and Vashist's convergence guarantees are derived for support vector machines and do not transfer formally to the network used here; the thesis adopts the paradigm, not its guarantees. The physiological signals, by contrast, are available at test time and are therefore an ordinary second modality rather than privileged information.


% ---------------------------------------------------------------------
% 5. Section 2.3 -- replace the paragraph that begins
%    "The physiological evidence available to such frameworks..."
% ---------------------------------------------------------------------
The physiological evidence available to such frameworks is also limited by existing datasets. Public multimodal stress datasets such as WESAD provide synchronized, high-quality wearable recordings with validated stress conditions, but they are small, demographically narrow, and based on stress induced in the laboratory~\cite{schmidt2018}. UBFC-Phys additionally records facial video from the same participants alongside wrist-worn BVP and EDA~\cite{sabour2023}, which makes it possible to study facial and physiological evidence together, but it shares these limitations, and this thesis uses an eight-participant pilot subset of it. Stress induced by the Trier Social Stress Test is a controlled proxy for, rather than a measure of, threat-related states, and results on so few participants generalize to new individuals only if each person's data are kept out of either training or testing.


% ---------------------------------------------------------------------
% 6. Section 2.1.10 and the 2.3 paragraph on the last three papers
% ---------------------------------------------------------------------
% The privileged-information sentences in 2.1.10 and in the 2.3 paragraph
% ("...using training-only scar annotations to enforce counterfactual invariance...")
% are now accurate for the final pipeline. Keep them, and delete any leftover note
% saying they apply "only if the mask is used during training alone".


% =====================================================================
% CHECKLIST: other Chapter 2 sentences that still describe the old
% CGF / WESAD+FFHQ pipeline. Update or delete each one.
% =====================================================================
% [ ] Kusner et al. relevance: "HRV and GSR" as the observable physiological signals
%     -> "BVP and EDA".
% [ ] Gohumpu et al. relevance and 2.1.3: "HRV and GSR as an efficient physiological
%     branch" -> "BVP and EDA from a wrist-worn device".
% [ ] Villarejo et al. relevance: "combines GSR with HRV and facial evidence"
%     -> "combines EDA with BVP and facial evidence".
% [ ] FairGRAPE relevance: the CGF Pruned30 DP results (0.0054 vs 0.0112 vs 0.0274)
%     belong to the obsolete CGF model -> delete those sentences, or restate them only
%     if EQUITAS-RCMF compression results exist.
% [ ] Ramesh et al. and Howard et al. relevance: references to "the compressed CGF model"
%     and "the thesis's compression experiments" -> describe the EQUITAS-RCMF deployment
%     (JIT compilation; INT8 quantization left as future work, Chapter 6).
% [ ] Lian and Celiktutan relevance: "point estimates on a single 80/20 split"
%     -> "point estimates from a single evaluation split".
% [ ] Liu et al. (JTT) relevance: "the thesis balances these groups directly during
%     training" (the old WeightedRandomSampler) -> check whether EQUITAS-RCMF training
%     balances the (scar, stress) groups; if not, delete that clause.
% [ ] Arjovsky et al. and Cohen et al. relevance: "the scar and threat groups are
%     balanced during construction" -> check against the UBFC-Phys pipeline, whose test
%     protocol varies the scar-label correlation (rho = 0.85 / 0.50 / 0.15).
% [ ] Anthis & Veitch and Arjovsky et al. relevance, and 2.3: sentences calling a
%     biased-train / inverted-test evaluation "future work" -> "the thesis defines such a
%     protocol (Chapter 3); its results will be reported in the final version".
% [ ] Any remaining use of the word "threat" as the label -> "stress", except where the
%     text discusses threat profiling as the motivating application.
